\section{Preparing the sparse-pair description}
\label{sec:predictor-setup}

We now prove Theorem~\ref{thm:predictor}. The proof is by induction
on the projective dimension \(d\). A large concentration on a proper
flat reduces the dimension. Otherwise, a small collection of flats
records all relevant concentrations, and the score in
Section~\ref{sec:predictor-score} gives the required description.

The following consequences of the parameter choices will be used
throughout:
\begin{equation}
\label{eq:parameter-orders}
 P=o(\sigma),\quad \log\sigma=o(P),\quad
 R\log L=o(P),\quad b+P\tau=o(L),\quad L=o(P).
\end{equation}
For example, \(P\le3D\sigma^{9\beta}\), whereas
\(D\le\sigma^{1-\eta/2}\); also \(b/L=O(\sigma^{-2\beta})\)
and \(P\tau/L=R\tau=o(1)\).
In particular, all fixed numerical constants used below are
eventually smaller than \(R\).

\subsection{Small sets and greedy removal of flats}

Orient the input so that \(n=|S|\le|T|\), and put \(n=q^{1+v}\).
Lemma~\ref{lem:incidence} gives \(n\le Cq^{(d+1)/2}\).
If \(n\le100qP\), send \(n\) and the rank of \(S\) among the
\(n\)-element subsets of the known set \(U_S\).
The elementary inequality
\(\binom{|U_S|}{n}\le(e|U_S|/n)^n\) bounds the message length by
\[
 O\bigl(\sigma+n(d_S+1)\bigr)
 \le CqP(d_S+P).
\]
The decoder takes \(W=S\). Henceforth \(n>100qP\), so \(v>0\).

Define
\[
 g_0=(v-\tfrac12)\sigma,\qquad g_1=(v-1)\sigma,\qquad
 \xi=L/100,\qquad \chi=P/100,\qquad \chi_0=P/10000.
\]
The following two greedy procedures refer to the original value
of \(n\), including in their thresholds.

\begin{enumerate}
\item If \(d=4\) and \(g_1>\chi_0\), repeatedly choose a hyperplane
containing the largest number of remaining points of \(S\).
Remove its remaining points whenever this number is at least \(q^2\).
Stop as soon as at least \(n/2\) points have been removed, or when
the maximum is below \(q^2\).
In the first case retain the removed points and their ordered list
of hyperplanes. In the second retain the residual and no list.
If the procedure is not used, retain all of \(S\).
Call the retained set \(S_1\).
\item If \(d\ge3\) and \(g_0>\chi_0\), apply the same rule to planes
and the points of \(S_1\), now with threshold
\[
 K_p=n^{4/3}q^{-1}\exp(-g_0/5).
\]
Stop on removing at least \(|S_1|/2\), or when the largest
remaining plane count is below \(K_p\).
Retain the removed points and the ordered plane list in the first
case, and the residual with no plane list in the second.
If this procedure is not used, retain \(S_1\).
Write \(S'\) for the final training set and \(n'=|S'|\).
\end{enumerate}

Ties are resolved by a fixed ordering of flats. Each retained list
partitions the training points by \emph{earliest membership}:
a point belongs to the first listed flat containing it. The
hyperplane partition, if present, is restricted to \(S'\) after
the second procedure. The sizes of the original hyperplane
cells, before that restriction, will also be used in the
geometric analysis; they are not needed by the decoder.

\begin{lemma}[Training data]
\label{lem:training}
The construction has \(n/4\le n'\le n\).
If the hyperplane list or plane list is present, its length satisfies
\[
 J_h\le n/q^2,\qquad
 J_p\le n/K_p=q^{1/2}\exp(-2g_0/15),
\]
respectively. If both lists are present, then
\(J_hJ_p\le q^{13/15+o(1)}\).
If a performed procedure retained a residual, that residual has
the corresponding flat cap.

If every training cell in both partitions has size less than
\(.04n'\), then for each query point \(x\) define \(O_x\) as the
union of the following cells: the cell of the earliest listed
hyperplane containing \(x\), and the cell of the earliest listed
plane containing \(x\). Omit a cell when no such flat exists.
Then \(|O_x|/n'<.08\). We call this the \emph{small-cell case}.

The ordered lists, \(n'\), all training cell sizes, and the table
of intersections between the two partitions can be communicated
in \(O(q)\) bits. From these data the decoder determines
\(|O_x|/n'\) for every \(x\), and determines which samples belong
to \(O_x\) from their flat memberships.
\end{lemma}

\begin{proof}
Every retention keeps at least half the previous set.
Each listed flat removes at least its threshold, and its assigned
points are disjoint from those of earlier flats, proving the
length bounds. In dimension four,
\[
 J_hJ_p
 \le q^{v-1}q^{1/2-(2/15)(v-1/2)}
 =q^{13v/15-13/30}
 \le q^{13/15+o(1)},
\]
because \(v\le3/2+O(1/\sigma)\).
Each single list has length at most \(q^{1/2+o(1)}\).
A flat is described by a basis in a fixed coordinate system, using
\(O(\sigma)\) bits; every cell or intersection size uses
\(O(\sigma)\) bits as well. These bounds give \(O(q)\) in total.
The union formula and the intersection table give \(|O_x|\).
For a sampled point, earliest membership in a known list determines
its cell. Finally the union involves at most two cells of size
less than \(.04n'\).
\end{proof}

\subsection{A large cell reduces the dimension}

Suppose a training cell has at least \(.04n'\) points.
Its listed flat \(K_0\), of projective dimension \(2\le j<d\),
satisfies
\[
 S_0=S\cap K_0,\qquad |S_0|\ge .01n.
\]
The encoder sends \(K_0\), with a canonical coordinate system.
The average incidence proportion from \(T\) into \(S_0\) is at
most \(100\tau/q\). Delete members whose proportion exceeds
\(\sqrt\tau/q\). This loses at most \(100\sqrt\tau\,|T|\) members,
so leaves at least \(|T|/2\) for large \(q\).
Each remaining covector restricts nontrivially to \(K_0\).

Normalize each nonzero projective restriction. Each restriction
has at most \(q^{d-j}\) projective lifts: after fixing one
nonzero restricted coordinate to be \(1\), the remaining \(d-j\)
coordinates are free. Among the dyads \([l,2l)\) of lift
multiplicity, choose one carrying at least \(|T|/(C\sigma)\)
of the retained lift mass. Let \(T_0\) be its set of distinct
restrictions. Then
\begin{equation}
\label{eq:reduced-pair}
 |S_0||T_0|\ge q^{j+1}\exp\{-b-O(\log\sigma)\},
 \qquad
 \frac{e(S_0,T_0)}{|S_0||T_0|}\le\frac{\sqrt\tau}{q}.
\end{equation}
Indeed, the number of restrictions is at least the retained mass
divided by \(2l\), and \(l\le q^{d-j}\).
The density assertion holds for each restriction separately,
so forgetting its multiplicity causes no problem.

The decoder's enclosures are
\[
 U_{S_0}=U_S\cap K_0,\qquad
 U_{T_0}=\{\text{nonzero projective restrictions with at least \(l\) lifts in \(U_T\)}\}.
\]
Their logarithmic gaps are bounded by
\(d_S+O(1)\) and \(d_T+O(\log\sigma)\).
The extra term in \(b\) in \eqref{eq:reduced-pair} is absorbed in
\(O(D\sigma^{6\beta})\). Moreover,
\[
 \sqrt\tau\le
 \sigma^{-200\cdot2^{j-2}\beta}\qquad(j<d).
\]
Thus the induction hypothesis applies after orienting the reduced
pair so that its smaller side is first. The dimension decreases,
and its initial value is at most four.

If the recursive output captures \(S_0\), it already captures
a constant fraction of \(S\) in a set of size at most \(ne^{CP}\).
Intersect it with \(U_S\), if necessary.
If instead it captures \(T_0\), let \(W_0\) be the decoded set and
\(C_0=W_0\cap T_0\) its hidden capture, with
\(|C_0|\ge c|T_0|\) and \(|W_0|/|C_0|\le e^{CP}\).
We convert this to a capture on the original side.

Take \(h=\lceil Cq(d_S+1)\rceil\) independent uniform samples
from \(C_0\), temporarily using the true hidden-set law.
On \(U_S\cap K_0\) define the cap
\[
 W=\{x:\text{fewer than } .2h/q
       \text{ of the sampled restrictions annihilate }x\}.
\]
By Lemma~\ref{lem:incidence}, at most \(C'q^{j+1}/|T_0|\)
points in the whole \(j\)-space have incidence proportion
less than \(.5/q\) into \(C_0\).
Every other point belongs to \(W\) with probability at most
\(\exp(-ch/q)\).
The bounds
\[
 |U_S\cap K_0|\le e^{d_S}n,\qquad
 n\le Cq^{j+1}/|T_0|
\]
therefore imply, by increasing \(C\) and applying Markov's
inequality, that
\[
 |W|\le C''q^{j+1}/|T_0|
 \le n\exp\{O(b+\log\sigma+1)\}
\]
with probability at least \(.9\).
The second displayed inequality for \(n\) follows by applying
mixing to \(S_0,T_0\) and using \(|S_0|\ge .01n\).

Every member of \(T_0\) has incidence proportion at most
\(\sqrt\tau/q\) into \(S_0\). Consequently every all-in-\(C_0\)
sample row, regardless of its bias, loses at most
\(5\sqrt\tau\,|S_0|\) points of \(S_0\) under this test:
sum the hit counts and divide by the rejection threshold.
Thus a row satisfying the size test also captures at least
\(.9|S_0|\) for large \(q\).

To implement this sampling with a message, send the integer \(h\)
and use a fresh table indexed by this row length.
Its \(O(\sigma)\)-bit cost is absorbed in the bounds below.
Propose independent rows uniform in the known \(W_0\).
A row falls entirely in \(C_0\)
and passes the size test with probability at least
\(.9e^{-ChP}\). Searching at most \(\exp(C'hP)\) rows gives
failure probability \(O(e^{-cq})\), and its row index takes
\(O(hP)\) bits. Conditional on being all in \(C_0\), a proposed
row is exactly iid uniform in \(C_0\); the comparison is made
before searching. After search, no unbiased-sampling assertion
is used.

The flat, dyad, recursive message, row length and optional conversion
have total length \(O(qP(d_S+d_T+P))\).
Here the recursive increases of \(O(\log\sigma)\) in the gaps
are absorbed because \(\log\sigma=o(P)\).
There are at most two dimension reductions, and each uses
fresh tables. The constants and failure probabilities can
therefore be absorbed in those of Theorem~\ref{thm:predictor}.

We have reduced the proof to training data satisfying the small-cell
case of Lemma~\ref{lem:training}. The next section bounds the rich
radial configurations that can correlate incidence tests at one
query point. Those bounds will be applied separately to validation
on \(S'\) and to output size in the ambient space.
